\chapter{2023年上海卷（春）}

\section{填空题}

\begin{problem}
已知集合 \(A=\{1,2\}\)，\(B=\{1,a\}\)，且 \(A=B\)，则 \(a=\)\fillinblank{} .
\end{problem}

\begin{answer}
\(2\)
\end{answer}

\begin{solution}
因为 \(A=B\)，所以两集合元素完全相同，又 \(1\in A\)，故 \(B\) 中已有元素 \(1\) 与之对应，由集合中元素互不相同得 \(a\ne 1\)，于是 \(a\) 只能等于 \(A\) 中另一元素 \(2\)，经检验 \(\{1,2\}=\{1,2\}\) 成立，所以 \(a=2\).
\end{solution}

\begin{problem}
已知向量 \(\overrightarrow{a}=(3,4)\)，\(\overrightarrow{b}=(1,2)\)，则 \(\overrightarrow{a}-2\overrightarrow{b}=\)\fillinblank{} .
\end{problem}

\begin{answer}
\(\left(1,0\right)\)
\end{answer}

\begin{solution}
由向量数乘与减法按坐标进行，得 \(\overrightarrow{a}-2\overrightarrow{b}=(3,4)-2(1,2)=(3-2,4-4)=\left(1,0\right)\).
\end{solution}

\begin{problem}
若不等式 \(\left|x-1\right|\le 2\)，则实数 \(x\) 的取值范围为\fillinblank{} .
\end{problem}

\begin{answer}
\(\left[-1,3\right]\)
\end{answer}

\begin{solution}
由 \(\left|x-1\right|\le 2\) 得 \(-2\le x-1\le 2\)，两边同时加 \(1\) 得 \(-1\le x\le 3\)，所以实数 \(x\) 的取值范围为 \(\left[-1,3\right]\).
\end{solution}

\begin{problem}
已知圆 \(C\) 的一般方程为 \(x^2+2x+y^2=0\)，则圆 \(C\) 的半径为\fillinblank{} .
\end{problem}

\begin{answer}
\(1\)
\end{answer}

\begin{solution}
对 \(x^2+2x+y^2=0\) 配方得 \(\left(x+1\right)^2-1+y^2=0\)，即 \(\left(x+1\right)^2+y^2=1\)，与圆的标准方程对照得圆心为 \(\left(-1,0\right)\)，半径为 \(1\).
\end{solution}

\begin{problem}
已知事件 \(A\) 发生的概率为 \(P(A)=0.5\)，则它的对立事件 \(\bar{A}\) 发生的概率 \(P(\bar{A})=\)\fillinblank{} .
\end{problem}

\begin{answer}
\(\frac{1}{2}\)
\end{answer}

\begin{solution}
因为 \(A\) 与 \(\bar{A}\) 为对立事件，所以 \(P(A)+P(\bar{A})=1\)，故 \(P(\bar{A})=1-P(A)=1-0.5=\frac{1}{2}\).
\end{solution}

\begin{problem}
已知正实数 \(a\)，\(b\) 满足 \(a+4b=1\)，则 \(ab\) 的最大值为\fillinblank{} .
\end{problem}

\begin{answer}
\(\frac{1}{16}\)
\end{answer}

\begin{solution}
因为 \(a>0\)，\(4b>0\)，由均值不等式得 \(1=a+4b\ge 2\sqrt{a\cdot 4b}=4\sqrt{ab}\)，即 \(ab\le \frac{1}{16}\)，当且仅当 \(a=4b\) 时取等号，结合 \(a+4b=1\) 得 \(a=\frac{1}{2}\)，\(b=\frac{1}{8}\) 时等号成立，所以 \(ab\) 的最大值为 \(\frac{1}{16}\).
\end{solution}

\begin{problem}
某校抽取 \(100\) 名学生测身高，其中身高最大值为 \(186\) cm，最小值为 \(154\) cm，根据身高数据绘制频率组距分布直方图，组距为 \(5\)，且第一组下限为 \(153.5\)，则组数为\fillinblank{} .
\end{problem}

\begin{answer}
\(7\)
\end{answer}

\begin{solution}
极差为 \(186-154=32\)，组距为 \(5\)，而 \(\frac{32}{5}=6.4\)，故至少需要 \(7\) 组，从 \(153.5\) 起依次为 \(\left[153.5+5(k-1),153.5+5k\right)\)，\(k=1,\cdots,6\)，末组为 \(\left[183.5,188.5\right]\)，最小值 \(154\) 落在第一组 \(\left[153.5,158.5\right)\) 内，最大值 \(186\) 落在第七组 \(\left[183.5,188.5\right]\) 内，恰覆盖全部数据，所以组数为 \(7\).
\end{solution}

\begin{problem}
设 \(\left(1-2x\right)^4=a_0+a_1x+a_2x^2+a_3x^3+a_4x^4\)，则 \(a_0+a_4=\)\fillinblank{} .
\end{problem}

\begin{answer}
\(17\)
\end{answer}

\begin{solution}
由二项式定理，\(\left(1-2x\right)^4\) 的展开通项为 \(\mathrm C_4^k1^{4-k}\left(-2x\right)^k\)，令 \(x=0\) 得常数项 \(a_0=1\)，取 \(k=4\) 得 \(a_4=\mathrm C_4^4\left(-2\right)^4=16\)，所以 \(a_0+a_4=1+16=17\).
\end{solution}

\begin{problem}
已知函数 \(f(x)=2^{-x}+1\)，且 \(g(x)= \begin{cases}
\log_2(x+1),  &x\ge 0,\\
f(-x),  &x<0,
\end{cases}\)
则方程 \(g(x)=2\) 的解为\fillinblank{} .
\end{problem}

\begin{answer}
\(x=3\)
\end{answer}

\begin{solution}
当\(x\ge 0\)时，\(g(x)=\log_2\left(x+1\right)\)，由\(g(x)=2\)得\(\log_2\left(x+1\right)=2\)，故\(x+1=4\)，\(x=3\)，满足\(x\ge 0\). 当\(x<0\)时，\(g(x)=f(-x)=2^{x}+1\)，由\(2^{x}+1=2\)得\(2^{x}=1\)，故\(x=0\)，舍去（不满足\(x<0\)）. 因此方程\(g(x)=2\)的解为\(x=3\).
\end{solution}

\begin{problem}
已知有 \(4\) 名男生 \(6\) 名女生，现从 \(10\) 人中任选 \(3\) 人，则恰有 \(1\) 名男生 \(2\) 名女生的概率为\fillinblank{} .
\end{problem}

\begin{answer}
\(\frac{1}{2}\)
\end{answer}

\begin{solution}
从\(10\)人中任选\(3\)人的等可能样本点数为\(\mathrm C_{10}^3=120\)，其中恰有\(1\)名男生\(2\)名女生的选法数为\(\mathrm C_4^1\mathrm C_6^2=4\times 15=60\)，故所求概率为\(\frac{60}{120}=\frac12\).
\end{solution}

\begin{problem}
设 \(z_1\)，\(z_2\in \mathbb{C}\) 且 \(z_1=\i\cdot z_2\)，满足 \(\left|z_1-1\right|=1\)，则 \(\left|z_1-z_2\right|\) 的取值范围为\fillinblank{} .
\end{problem}

\begin{answer}
按本页现稿 \(\left[0,2\sqrt2\right]\)（原PDF题面为 \(\left[0,2+\sqrt2\right]\)）.
\end{answer}

\begin{solution}
设\(z_1=a+b\i\)，\(z_2=c+d\i\)，其中\(a\)，\(b\)，\(c\)，\(d\)为实数. 按本页现稿\(z_1=\i z_2\)，得\(a+b\i=\i\left(c+d\i\right)=-d+c\i\)，故\(a=-d\)，\(b=c\)，即\(z_2=b-a\i\). 由\(\left|z_1-1\right|=1\)得\(\left(a-1\right)^2+b^2=1\)，即\(a^2+b^2=2a\)，且\(a\in\left[0,2\right]\). 又\(\left|z_1-z_2\right|^2=\left(a-b\right)^2+\left(a+b\right)^2=2\left(a^2+b^2\right)=4a\)，故\(\left|z_1-z_2\right|=2\sqrt{a}\in\left[0,2\sqrt2\right]\)，两端分别在\(z_1=0\)，\(z_2=0\)与\(z_1=2\)，\(z_2=-2\i\)处取到. 原PDF题面为\(z_1=\i\cdot\overline{z_2}\)，此时\(a=d\)，\(b=c\)，\(z_2=b+a\i\)，\(\left|z_1-z_2\right|=\sqrt{2}\left|a-b\right|\)，同由\(\left(a-1\right)^2+b^2=1\)得取值范围为\(\left[0,2+\sqrt2\right]\). 按本页现稿结论为\(\left[0,2\sqrt2\right]\).
\end{solution}

\begin{problem}
已知空间向量 \(\overrightarrow{OA}\)，\(\overrightarrow{OB}\)，\(\overrightarrow{OC}\) 都是单位向量，且 \(\overrightarrow{OA}\perp \overrightarrow{OB}\)，\(\overrightarrow{OA}\perp \overrightarrow{OC}\)，\(\overrightarrow{OB}\) 与 \(\overrightarrow{OC}\) 的夹角为 \(60^\circ\). 若 \(P\) 为空间任意一点，且 \(\left|\overrightarrow{OP}\right|=1\)，满足 \(\left|\overrightarrow{OP}\cdot\overrightarrow{OC}\right|\le \left|\overrightarrow{OP}\cdot\overrightarrow{OB}\right|\le \left|\overrightarrow{OP}\cdot\overrightarrow{OA}\right|\)，则 \(\overrightarrow{OP}\cdot\overrightarrow{OC}\) 的最大值为\fillinblank{} .
\end{problem}

\begin{answer}
\(\frac{\sqrt{21}}{7}\)
\end{answer}

\begin{solution}
因为\(\overrightarrow{OA}\perp\overrightarrow{OB}\)，\(\overrightarrow{OA}\perp\overrightarrow{OC}\)，所以\(\overrightarrow{OA}\)垂直于平面\(BOC\). 以\(O\)为原点，\(\overrightarrow{OA}\)为\(z\)轴，\(\overrightarrow{OB}\)为\(y\)轴建立右手直角坐标系，则\(\overrightarrow{OA}=(0,0,1)\)，\(\overrightarrow{OB}=(0,1,0)\)，\(\overrightarrow{OC}\)与\(z\)轴垂直且为单位向量、与\(\overrightarrow{OB}\)夹角为\(60^\circ\)，故经\(x\mapsto -x\)镜像可设\(\overrightarrow{OC}=\left(\frac{\sqrt{3}}{2},\frac12,0\right)\). 设\(\overrightarrow{OP}=(x,y,z)\)，\(x^2+y^2+z^2=1\)，记\(w=\overrightarrow{OP}\cdot\overrightarrow{OC}=\frac{\sqrt{3}}{2}x+\frac12y\)，约束为\(\left|w\right|\le \left|y\right|\le \left|z\right|\). 由\(y^2\le z^2=1-x^2-y^2\)得\(x^2+2y^2\le 1\). 若\(w\le 0\)则\(w<\frac{\sqrt{21}}{7}\)，下设\(w>0\)，记\(s=\sqrt{3}x+y=2w>0\)，则\(\left|y\right|\ge \frac{s}{2}\)，且由\(\sqrt{3}x=s-y\)得\(\frac{\left(s-y\right)^2}{3}+2y^2\le 1\)，即\(7y^2-2sy+s^2-3\le 0\). 记\(f(y)=7y^2-2sy+s^2-3\)，其顶点在\(y=\frac{s}{7}\). 当\(y\ge \frac{s}{2}\)时\(f\)单调递增，故\(f(y)\ge f\left(\frac{s}{2}\right)=\frac{7s^2}{4}-3\)，由\(f(y)\le 0\)得\(s^2\le \frac{12}{7}\). 当\(y\le -\frac{s}{2}\)时\(f(y)\ge f\left(-\frac{s}{2}\right)=\frac{15s^2}{4}-3\)，得\(s^2\le \frac45<\frac{12}{7}\). 于是\(s\le \frac{2\sqrt{21}}{7}\)，即\(w\le \frac{\sqrt{21}}{7}\). 取\(P\left(\frac{\sqrt{7}}{7},\frac{\sqrt{21}}{7},\frac{\sqrt{21}}{7}\right)\)，则\(\left|\overrightarrow{OP}\right|=1\)，\(\overrightarrow{OP}\cdot\overrightarrow{OC}=\left|\overrightarrow{OP}\cdot\overrightarrow{OB}\right|=\left|\overrightarrow{OP}\cdot\overrightarrow{OA}\right|=\frac{\sqrt{21}}{7}\)，满足约束且取到上界. 故最大值为\(\frac{\sqrt{21}}{7}\).
\end{solution}
\section{单选题}

\begin{problem}
下列函数是偶函数的是（ \quad ）.

\choices
    {\(y=\sin x\)}
    {\(y=\cos x\)}
    {\(y=x^3\)}
    {\(y=2^x\)}
\end{problem}

\begin{answer}
B
\end{answer}

\begin{solution}
偶函数满足定义域关于原点对称且\(f(-x)=f(x)\). \(y=\sin x\)与\(y=x^3\)满足\(f(-x)=-f(x)\)，为奇函数；\(y=\cos x\)满足\(\cos\left(-x\right)=\cos x\)，为偶函数；\(y=2^x\)在\(x=1\)处\(2^{-1}\ne 2^1\)且\(2^{-1}\ne -2^1\)，为非奇非偶函数. 故选B.
\end{solution}

\begin{problem}
如图为 \(2018-2021\) 年中国货物进出口总额的条形统计图，则下列对于进出口贸易额描述错误的是（ \quad ）.

\texfigure[max width=0.48\linewidth]{img_repaint/2023/shanghai_spring/q14_fig1.tex}

\choices
    {从 \(2018\) 年开始，\(2021\) 年的进出口总额增长率最大}
    {从 \(2018\) 年开始，进出口总额逐年增大}
    {从 \(2018\) 年开始，进口总额逐年增大}
    {从 \(2018\) 年开始，\(2020\) 年的进出口总额增长率最小}
\end{problem}

\begin{answer}
C
\end{answer}

\begin{solution}
原PDF条形图中浅色段为进口，深色段为出口，各年进口额依次为\(14.09\)，\(14.33\)，\(14.29\)，\(17.37\)（万亿），出口额依次为\(16.41\)，\(17.24\)，\(17.93\)，\(21.73\)（万亿），进出口总额依次为\(30.50\)，\(31.57\)，\(32.22\)，\(39.10\)（万亿），逐年增大，故B正确. 进口额在\(2020\)年由\(14.33\)降为\(14.29\)，并非逐年增大，故C的描述错误. 以\(2018\)年为基期，\(2019\)年、\(2020\)年、\(2021\)年总额增长率分别约为\(3.51\%\)，\(2.06\%\)，\(21.35\%\)，故\(2021\)年最大、\(2020\)年最小，A与D正确. 因此描述错误的是C.
\end{solution}

\begin{problem}
如图所示，在正方体 \(ABCD-A_1B_1C_1D_1\) 中，点 \(P\) 为线段 \(A_1C_1\) 上的动点，则下列直线中，始终与直线 \(BP\) 异面的是（ \quad ）.

\bitmapfigure[max width=0.38\linewidth]{img/2023/shanghai_spring/q15_fig1.png}

\choices
    {\(DD_1\)}
    {\(AC\)}
    {\(AD_1\)}
    {\(B_1C\)}
\end{problem}

\begin{answer}
B
\end{answer}

\begin{solution}
记底面为\(ABCD\)，顶面为\(A_1B_1C_1D_1\)且\(AA_1\)垂直于底面. 直线\(AC\)在平面\(ACC_1A_1\)内，点\(B\)不在该平面内，而动点\(P\)恒在\(A_1C_1\)上，\(A_1C_1\)在平面\(ACC_1A_1\)内且\(P\notin AC\)，故直线\(BP\)与平面\(ACC_1A_1\)交于\(P\)，交点不在\(AC\)上，所以\(BP\)与\(AC\)恒异面. 当\(P\)为\(A_1C_1\)中点（即\(B_1D_1\)中点）时，\(BP\)与\(DD_1\)同在平面\(BDD_1B_1\)内；当\(P\)与\(C_1\)重合时，\(BP\)即\(BC_1\)，与\(AD_1\)平行（由\(AB\parallel D_1C_1\)且\(AB=D_1C_1\)得四边形\(ABC_1D_1\)为平行四边形），且与\(B_1C\)为同一面\(BCC_1B_1\)的两条对角线而相交. 故\(DD_1\)，\(AD_1\)，\(B_1C\)均非始终异面. 因此始终与\(BP\)异面的是\(AC\)，选B.
\end{solution}

\begin{problem}
已知数列 \(\{a_n\}\) 的各项均为实数，\(S_n\) 为其前 \(n\) 项和，若对任意 \(k>2022\)，都有 \(\left|S_k\right|>\left|S_{k+1}\right|\)，则下列说法正确的是（ \quad ）.

\choices
    {\(a_1,a_3,a_5,\cdots,a_{2n-1}\) 为等差数列，\(a_2,a_4,a_6,\cdots,a_{2n}\) 为等比数列}
    {\(a_1\)，\(a_3\)，\(a_5\)，\(\cdots\)，\(a_{2n-1}\) 为等比数列，\(a_2\)，\(a_4\)，\(a_6\)，\(\cdots\)，\(a_{2n}\) 为等差数列}
    {\(a_1\)，\(a_2\)，\(a_3\)，\(\cdots\)，\(a_{2022}\) 为等差数列，\(a_{2022}\)，\(a_{2023}\)，\(a_{2024}\)，\(\cdots\)，\(a_n\) 为等比数列}
    {\(a_1\)，\(a_2\)，\(a_3\)，\(\cdots\)，\(a_{2022}\) 为等比数列，\(a_{2022}\)，\(a_{2023}\)，\(a_{2024}\)，\(\cdots\)，\(a_n\) 为等差数列}
\end{problem}

\begin{answer}
C
\end{answer}

\begin{solution}
以下按“四项中唯一能与条件相容者为正确”的出题意图证明. 条件要求\(\left|S_k\right|>\left|S_{k+1}\right|\)对一切\(k>2022\)成立，故\(\left(\left|S_k\right|\right)_{k>2022}\)严格递减且有下界\(0\)，特别\(\left(S_k\right)_{k>2022}\)有界，且\(a_k\ne 0\)对一切\(k>2023\)成立，否则将出现相邻两项绝对值相等.

对A，记奇子列\(c_j=a_{2j-1}=c_1+\left(j-1\right)d\)，偶子列为等比数列（首项\(A\)，公比\(q\)），则\(S_{2m}=O_m+E_m\)，其中\(O_m=mc_1+\frac{m\left(m-1\right)}{2}d\)，\(q\ne 1\)时\(E_m=\frac{A\left(1-q^m\right)}{1-q}\)，\(q=1\)时\(E_m=mA\). 由\(a_k\ne 0\)，子列全零的情形已不可能. 若\(d\ne 0\)，则\(\left|O_m\right|\)二次无界：\(\left|q\right|<1\)、\(q=-1\)、\(q=0\)时\(E_m\)有界，\(q=1\)时\(E_m\)线性，二次项占优使\(\left|S_{2m}\right|\)无界；\(\left|q\right|>1\)时\(E_m=\alpha+\beta q^m\)（\(A=0\)则\(E_m=0\)，同样无界），\(\left|S_{2m}\right|\ge \left|\beta\right|\left|q\right|^m-\left|O_m\right|-\left|\alpha\right|\to\infty\). 若\(d=0\)，则奇子列为非零常数\(c\)，\(O_m=mc\)线性无界：\(E_m\)有界或指数增长时\(\left|S_{2m}\right|\)无界；\(q=1\)时\(S_{2m}=\left(c+A\right)m\)，除非\(A=-c\)同样无界，而\(A=-c\)时\(S_{2m}\)与\(S_{2m+1}\)均恒定，绝对值不再严格递减. 故A不可能，同理偶子列为等差的B不可能. 对D，记尾段\(b_m=a_{2021+m}=b_1+\left(m-1\right)d\)，则\(S_{2021+m}=S_{2021}+mb_1+\frac{m\left(m-1\right)}{2}d\)，\(d\)与\(b_1\)不全为零时\(\left|S_n\right|\)无界，全为零时尾段出现相邻绝对值相等，故D不可能.

对C，取尾段\(a_{2022}=-1\)，\(a_{2023}=-\frac12\)，\(a_{2024}=-\frac14\)，\(\cdots\)（首项为\(-1\)、公比为\(\frac12\)的等比数列），并取\(S_{2021}=10\)，则\(S_{2021+m}=8+2^{1-m}\)（\(m\ge 1\)），即\(S_{2022}=9\)，\(S_{2023}=8.5\)，\(S_{2024}=8.25\)，\(\cdots\)为正数严格递减；前段取\(a_1=\frac{340}{337}\)，\(a_{2022}=-1\)，和为\(9\)的等差数列（公差为\(-\frac{677}{681077}\)），恰与\(S_{2021}=10\)衔接. 此时条件对一切\(k>2022\)成立. 故四项中唯一可与条件相容的是C，选C.
\end{solution}
\section{解答题}

\begin{problem}
已知三棱锥 \(P-ABC\) 中，\(AB\perp AC\)，\(PA\perp\) 平面 \(ABC\)，\(PA=AB=3\)，\(AC=4\)，\(M\) 为 \(BC\) 中点，过点 \(M\) 分别作平行于平面 \(PAB\) 的直线交 \(AC\)，\(PC\) 于点 \(E\)，\(F\).

\begin{enumerate}
\item 求直线 \(PM\) 与平面 \(ABC\) 所成的角的正切值；
\item 证明：平面 \(MEF\parallel\) 平面 \(PAB\)，并求直线 \(ME\) 到平面 \(PAB\) 的距离.
\end{enumerate}

\FigureLayoutDeclare{img/2023/shanghai_spring/q17_fig1.png}{5.0cm}{20bp 31bp 60bp 5bp}
\bitmapfigure[max width=0.42\linewidth]{img/2023/shanghai_spring/q17_fig1.png}
\end{problem}

\begin{solution}
\begin{enumerate}
\item 因为 \(PA\perp\) 平面 \(ABC\)，所以 \(AM\) 为 \(PM\) 在平面 \(ABC\) 内的射影，故 \(\angle PMA\) 即为直线 \(PM\) 与平面 \(ABC\) 所成的角. 因为 \(AB\perp AC\)，\(AB=3\)，\(AC=4\)，所以 \(BC=5\)，而 \(M\) 为 \(BC\) 中点，故 \(AM=\frac{5}{2}\). 在 \(\mathrm{Rt}\triangle PAM\) 中，\(\tan\angle PMA=\frac{PA}{AM}=\frac{6}{5}\). 所以直线 \(PM\) 与平面 \(ABC\) 所成的角的正切值为 \(\frac{6}{5}\).
\item 设 \(\Pi\) 为过点 \(M\) 且平行于平面 \(PAB\) 的平面，则 \(\Pi\) 唯一确定. 平面 \(\Pi\) 与平面 \(ABC\) 都过点 \(M\)，其交线过点 \(M\) 且平行于两平行平面的交线 \(AB\)，该交线交线段 \(AC\) 于一点，记为 \(E\)，则 \(ME\parallel AB\)，由中位线定理得 \(E\) 为 \(AC\) 中点. 同理，平面 \(\Pi\) 与平面 \(PBC\) 的交线过点 \(M\) 且平行于交线 \(PB\)，交线段 \(PC\) 于一点，记为 \(F\)，则 \(MF\parallel PB\)，\(F\) 为 \(PC\) 中点. 于是 \(ME\)，\(MF\) 都平行于平面 \(PAB\)，且 \(ME\cap MF=M\)，故平面 \(MEF\) 即为 \(\Pi\)，所以平面 \(MEF\parallel\) 平面 \(PAB\). 因为 \(PA\perp\) 平面 \(ABC\)，所以 \(PA\perp AC\)，又 \(AB\perp AC\)，且 \(PA\)，\(AB\subset\) 平面 \(PAB\)，\(PA\cap AB=A\)，故 \(AC\perp\) 平面 \(PAB\). 由 \(ME\parallel\) 平面 \(PAB\) 可知直线 \(ME\) 到平面 \(PAB\) 的距离等于点 \(E\) 到平面 \(PAB\) 的距离，而 \(E\in AC\) 且 \(AC\perp\) 平面 \(PAB\) 于点 \(A\)，故所求距离为 \(EA=\frac{1}{2}AC=2\).
\end{enumerate}
\end{solution}

\begin{problem}
在 \(\triangle ABC\) 中，内角 \(A\)，\(B\)，\(C\) 所对的边为 \(a\)，\(b\)，\(c\)，其中 \(b=2\).

\begin{enumerate}
\item 若 \(A+C=60^\circ\)，且 \(a=2c\)，求边长 \(c\) 的值；
\item 若 \(A-C=30^\circ\)，\(a=\sqrt2\,c\sin A\)，求 \(S_{\triangle ABC}\).
\end{enumerate}
\end{problem}

\begin{solution}
\begin{enumerate}
\item 因为 \(A+C=60^\circ\)，所以 \(B=120^\circ\). 由余弦定理 \(\cos B=\frac{a^2+c^2-b^2}{2ac}\)，代入 \(a=2c\)，\(b=2\) 得 \(\frac{5c^2-4}{4c^2}=-\frac{1}{2}\)，解得 \(c^2=\frac{4}{7}\). 由 \(c>0\) 得 \(c=\frac{2\sqrt7}{7}\).
\item 由 \(a=\sqrt2\,c\sin A\) 及正弦定理 \(\frac{a}{\sin A}=\frac{c}{\sin C}\) 得 \(\sin C=\frac{1}{\sqrt2}\). 故 \(C=45^\circ\) 或 \(C=135^\circ\)，若 \(C=135^\circ\)，则 \(A=C+30^\circ=165^\circ\)，\(A+C>180^\circ\)，不可能. 所以 \(C=45^\circ\)，进而 \(A=75^\circ\)，\(B=60^\circ\). 由正弦定理 \(c=\frac{b\sin C}{\sin B}=\frac{2\sqrt6}{3}\)，且 \(\sin75^\circ=\sin45^\circ\cos30^\circ+\cos45^\circ\sin30^\circ=\frac{\sqrt6+\sqrt2}{4}\). 故 \(S_{\triangle ABC}=\frac{1}{2}bc\sin A=\frac{1}{2}\cdot 2\cdot\frac{2\sqrt6}{3}\cdot\frac{\sqrt6+\sqrt2}{4}=\frac{\sqrt3+3}{3}\).
\end{enumerate}
\end{solution}

\begin{problem}
已知 \(S\) 为正比例系数，定义 \(S=\frac{F_0}{V_0}\)，\(F_0\) 为建筑物暴露在空气中的面积（单位：平方米），\(V_0\) 为建筑物的体积（单位：立方米）.

\begin{enumerate}
\item 若有一个圆柱体建筑的底面半径为 \(R\)，高度为 \(H\)，求该建筑体的 \(S\)（用 \(R\)，\(H\) 表示）；
\item 现有一个建筑体，侧面皆垂直于地面，设 A 为底面面积，\(L\) 为建筑底面周长. 已知 \(f\) 为正比例系数，\(L^2\) 与 A 成正比，定义 \(f=\frac{L^2}{A}\)，建筑面积即为每一层的底面面积，总建筑面积即为每层建筑面积之和，值为 \(T\). 已知该建筑体推导得出 \(S=\sqrt{\frac{f\cdot n}{T}+\frac{1}{3n}}\)，\(n\) 为层数，层高为 \(3\) 米，其中 \(f=18\)，\(T=10000\)，试求当取第几层时，该建筑体 \(S\) 最小？
\end{enumerate}
\end{problem}

\begin{solution}
\begin{enumerate}
\item 圆柱侧面积为 \(2\pi RH\)，暴露在空气中的表面还有上底面 \(\pi R^2\)（下底面不计），故 \(F_0=2\pi RH+\pi R^2\)，\(V_0=\pi R^2H\)，所以 \(S=\frac{F_0}{V_0}=\frac{2H+R}{HR}\).
\item 按原PDF题面，\(S=\sqrt{\frac{fn}{T}}+\frac{1}{3n}\)，代入 \(f=18\)，\(T=10000\) 得 \(\psi\left(n\right)=\frac{3\sqrt{2n}}{100}+\frac{1}{3n}\)，其中 \(n\) 为正整数. 考虑 \(\psi\) 在 \(\left(0,+\infty\right)\) 上的连续延拓，\(\psi'\left(x\right)=\frac{3\sqrt2}{200\sqrt x}-\frac{1}{3x^2}\)，令导数为零得 \(x_0=\left(\frac{100\sqrt2}{9}\right)^{\frac{2}{3}}\approx6.274\)，\(\psi'\) 在 \(\left(0,x_0\right)\) 上为负，在 \(\left(x_0,+\infty\right)\) 上为正，故 \(\psi\) 先减后增. 正整数 \(n\) 的最小值在 \(6\) 或 \(7\) 处取得，\(\psi\left(6\right)=\frac{3\sqrt3}{50}+\frac{1}{18}\approx0.159\)，\(\psi\left(7\right)=\frac{3\sqrt{14}}{100}+\frac{1}{21}\approx0.160\)，\(\psi\left(6\right)<\psi\left(7\right)\). 所以当取 \(6\) 层时，该建筑体 \(S\) 最小. 注：若按本页现稿 \(S=\sqrt{\frac{fn}{T}+\frac{1}{3n}}\)，则最小化 \(S\) 等价于最小化 \(\varphi\left(n\right)=\frac{18n}{10000}+\frac{1}{3n}\)，\(\varphi\) 在约 \(13.6\) 处由减转增，比较整数点得 \(\varphi\left(14\right)\) 最小，即按现稿公式应在 \(14\) 层时 \(S\) 最小.
\end{enumerate}
\end{solution}

\begin{problem}
已知椭圆 \(\Gamma:\frac{x^2}{m^2}+\frac{y^2}{3}=1\) （\(m>0\)，\(\ m\neq \sqrt3\)）.

\begin{enumerate}
\item 若 \(m=2\)，求椭圆 \(\Gamma\) 的离心率；
\item 设 \(A_1\)，\(A_2\) 为椭圆 \(\Gamma\) 的左右顶点，若椭圆 \(\Gamma\) 上一点 \(E\) 的纵坐标为 \(1\)，且 \(\overrightarrow{EA_1}\cdot\overrightarrow{EA_2}=-2\)，求 \(m\) 的值；
\item 若 \(P\) 为椭圆 \(\Gamma\) 上一点，过点 \(P\) 作一条斜率为 \(\sqrt3\) 的直线与双曲线 \(\frac{y^2}{5m^2}-\frac{x^2}{5}=1\) 仅有一个公共点，求 \(m\) 的取值范围.
\end{enumerate}
\end{problem}

\begin{solution}
\begin{enumerate}
\item 当 \(m=2\) 时，椭圆方程为 \(\frac{x^2}{4}+\frac{y^2}{3}=1\)，故 \(a=2\)，\(b=\sqrt3\)，\(c=\sqrt{a^2-b^2}=1\)，离心率 \(e=\frac{c}{a}=\frac{1}{2}\).
\item 记 \(A_1\left(-m,0\right)\)，\(A_2\left(m,0\right)\). 由 \(\frac{x^2}{m^2}+\frac{1}{3}=1\) 得点 \(E\) 的横坐标 \(x=\pm\frac{\sqrt6}{3}m\). 于是 \(\overrightarrow{EA_1}\cdot\overrightarrow{EA_2}=x^2-m^2+1=1-\frac{m^2}{3}\). 令其等于 \(-2\) 得 \(m^2=9\)，由 \(m>0\) 得 \(m=3\)（满足 \(m\neq\sqrt3\)）.
\item 设过点 \(P\) 的直线为 \(l\)：\(y=\sqrt3x+b\)，双曲线的渐近线为 \(y=\pm mx\). 若 \(l\) 与渐近线平行，则 \(\sqrt3=m\)，与 \(m\neq\sqrt3\) 矛盾，故 \(l\) 与渐近线不平行，此时 \(3-m^2\neq0\). 联立 \(l\) 与双曲线得 \(\left(3-m^2\right)x^2+2\sqrt3bx+b^2-5m^2=0\)，仅有一个公共点当且仅当判别式为零，化简得 \(b^2=5m^2-15\geqslant0\)，故 \(m>\sqrt3\). 又 \(P\) 在椭圆上，故 \(l\) 与椭圆有公共点，联立得 \(\left(3+3m^2\right)x^2+2\sqrt3bm^2x+\left(b^2-3\right)m^2=0\) 有实数解，其判别式非负化简得 \(b^2\leqslant3m^2+3\). 于是 \(5m^2-15\leqslant3m^2+3\)，得 \(m\leqslant3\). 对每个 \(m\in\left(\sqrt3,3\right]\)，取 \(b^2=5m^2-15\) 即同时满足与双曲线相切及与椭圆有交点. 故 \(m\) 的取值范围为 \(\left(\sqrt3,3\right]\).
\end{enumerate}
\end{solution}

\begin{problem}
设函数 \(f(x)=ax^3-(a+1)x^2+x\)，\(g(x)=kx+m\)，其中 \(a\ge 0\)，\(k\)，\(m\in \R\)，若任意 \(x\in[0,1]\) 均有 \(f(x)\le g(x)\)，则称函数 \(y=g(x)\) 是函数 \(y=f(x)\) 的控制函数，且对于所有满足条件的函数 \(y=g(x)\) 在 \(x\) 处取得的最小值记为 \(\bar f(x)\).

\begin{enumerate}
\item 若 \(a=2\)，\(g(x)=x\)，试问 \(y=g(x)\) 是否为 \(y=f(x)\) 的控制函数；
\item 若 \(a=0\)，使得直线 \(y=h(x)\) 是曲线 \(y=f(x)\) 在 \(x=\frac14\) 处的切线，证明：函数 \(y=h(x)\) 为函数 \(y=f(x)\) 的控制函数，并求 \(\bar f\left(\frac14\right)\) 的值；
\item 若曲线 \(y=f(x)\) 在 \(x=x_0\)（\(x_0\in(0,1)\)）处的切线过点 \((1,0)\)，且 \(c\in[x_0,1]\)，证明：当且仅当 \(c=x_0\) 或 \(c=1\) 时，\(\bar f(c)=f(c)\).
\end{enumerate}
\end{problem}

\begin{solution}
\begin{enumerate}
\item 当 \(a=2\)，\(g\left(x\right)=x\) 时，\(f\left(x\right)-g\left(x\right)=2x^3-3x^2\)，其导数为 \(6x\left(x-1\right)\leqslant0\)（\(x\in\left[0,1\right]\)），故 \(f-g\) 在 \(\left[0,1\right]\) 上单调递减，又在 \(x=0\) 处取值为 \(0\)，所以 \(f\left(x\right)\leqslant g\left(x\right)\) 在 \(\left[0,1\right]\) 上恒成立. 故 \(y=g\left(x\right)\) 是 \(y=f\left(x\right)\) 的控制函数.
\item 当 \(a=0\) 时，\(f\left(x\right)=-x^2+x\)，\(f'\left(x\right)=-2x+1\)，\(f\left(\frac{1}{4}\right)=\frac{3}{16}\)，\(f'\left(\frac{1}{4}\right)=\frac{1}{2}\)，切线为 \(h\left(x\right)=\frac{1}{2}x+\frac{1}{16}\). 记 \(\psi\left(x\right)=f\left(x\right)-h\left(x\right)=-\left(x-\frac{1}{4}\right)^2\leqslant0\)，等号仅在 \(x=\frac{1}{4}\) 处成立，故 \(h\) 是控制函数. 对任意控制函数 \(g\) 有 \(g\left(\frac{1}{4}\right)\geqslant f\left(\frac{1}{4}\right)\)，而 \(h\left(\frac{1}{4}\right)=f\left(\frac{1}{4}\right)\)，故下确界可以取到，\(\bar f\left(\frac{1}{4}\right)=f\left(\frac{1}{4}\right)=\frac{3}{16}\).
\item 由 \(f\left(x\right)-f\left(1\right)=\left(x-1\right)x\left(ax-1\right)\) 及切线过点 \(\left(1,0\right)\) 得 \(f'\left(x_0\right)=x_0\left(ax_0-1\right)\)，即 \(3ax_0^2-2\left(a+1\right)x_0+1=ax_0^2-x_0\)，化简得 \(\left(2ax_0-1\right)\left(x_0-1\right)=0\). 由 \(x_0\neq1\) 得 \(a=\frac{1}{2x_0}>\frac{1}{2}\)，\(x_0=\frac{1}{2a}\)，切线斜率 \(f'\left(x_0\right)=-\frac{1}{4a}\)，切线为 \(t\left(x\right)=-\frac{1}{4a}\left(x-1\right)\). 直接计算得 \(t\left(x\right)-f\left(x\right)=-a\left(x-1\right)\left(x-x_0\right)^2\geqslant0\)（\(x\in\left[0,1\right]\)），等号仅在 \(x=x_0\) 或 \(x=1\) 处成立，故 \(t\) 是控制函数. 设 \(g\left(x\right)=kx+m\) 为任一控制函数，记 \(h=g-t\)，则 \(h\) 为一次函数（或常数），且 \(h\left(x_0\right)\geqslant0\)，\(h\left(1\right)\geqslant0\). 对 \(c\in\left[x_0,1\right]\)，\(h\left(c\right)\) 为端点值的凸组合，故 \(h\left(c\right)\geqslant0\)，即 \(g\left(c\right)\geqslant t\left(c\right)\). 当 \(c\in\left(x_0,1\right)\) 时 \(t\left(c\right)>f\left(c\right)\)，故一切控制函数在 \(c\) 处取值都严格大于 \(f\left(c\right)\)，\(\bar f\left(c\right)\geqslant t\left(c\right)>f\left(c\right)\). 而在 \(c=x_0\) 与 \(c=1\) 处 \(t\) 与 \(f\) 相合，故 \(\bar f=f\). 于是对 \(c\in\left[x_0,1\right]\)，当且仅当 \(c=x_0\) 或 \(c=1\) 时 \(\bar f\left(c\right)=f\left(c\right)\).
\end{enumerate}
\end{solution}
